%Two resources useful for abstract writing.
% Guidance of how to write an abstract/summary provided by Nature: https://cbs.umn.edu/sites/cbs.umn.edu/files/public/downloads/Annotated_Nature_abstract.pdf %https://writingcenter.gmu.edu/guides/writing-an-abstract
\chapter*{\center \Large  Abstrait}
%%%%%%%%%%%%%%%%%%%%%%%%%%%%%%%%%%%%%%
% Replace all text with your text
%%%%%%%%%%%%%%%%%%%%%%%%%%%%%%%%%%%
Ce rapport présente la conception et le développement d'une application web suivant une architecture 3-tiers. L'objectif principal de ce projet est de mettre en pratique les concepts fondamentaux des architectures logicielles distribuées, en séparant clairement les couches de présentation, de logique métier et de données.

L'application est construite autour de technologies modernes et largement utilisées dans l'industrie. La couche de présentation utilise JavaScript pour offrir une interface utilisateur interactive et réactive. La couche métier est implémentée en Python, assurant le traitement des requêtes et l'application des règles fonctionnelles. Enfin, MongoDB, une base de données NoSQL orientée documents, est employée pour la persistance des données, offrant flexibilité et scalabilité.

Ce projet met en évidence les avantages de l'architecture 3-tiers, notamment la modularité, la maintenabilité et la possibilité de faire évoluer chaque couche indépendamment. Nous détaillons les choix techniques effectués, les défis rencontrés lors de l'implémentation, ainsi que les solutions adoptées pour garantir la cohérence et la performance du système.

Les résultats obtenus démontrent l'efficacité de cette approche architecturale pour le développement d'applications web robustes et évolutives, tout en respectant les bonnes pratiques de génie logiciel.
%%%%%%%%%%%%%%%%%%%%%%%%%%%%%%%%%%%%%%%%%%%%%%%%%%%%%%%%%%%%%%%%%%%%%%%%%s
~\\[1cm]
\noindent % Provide your key words
\textbf{Mots-clés:} architecture 3-tiers, base de données, JavaScript, Python, MongoDB

\vfill
\noindent
